\documentclass[10pt,a4paper]{amsart}
\usepackage[margin=19mm]{geometry}
\usepackage{amsmath,amssymb,mathtools}
\usepackage[scr=rsfs,cal=euler]{mathalfa}
\usepackage{xfrac}
\usepackage[unicode,hidelinks,bookmarks=false]{hyperref}
\newcommand{\ie}{i.e.\ }
\newcommand{\cf}{cf.\ }
\newcommand{\resp}{respectively\ }
\newcommand{\Subsection}{\subsection}
\providecommand{\nobreakdash}{\nobreak\hbox{-}}
\setlength{\emergencystretch}{2em}
\multlinegap0pt
\usepackage{amssymb}
\usepackage{braket}
\usepackage{float}
\newtheorem{lemma}{Lemma}
\newcommand{\Tr}{\mathrm{Tr}}
\title{\LARGE \bf Noise Resilience and Robust Convergence Guarantees for the Variational Quantum Eigensolver}
\author{Mirko Legnini, Julian Berberich}
\date{Source: arXiv:2601.16758v2 (CC BY 4.0)}
\begin{document}
\maketitle

\noindent\emph{Source and licence.} \LARGE \bf Noise Resilience and Robust Convergence Guarantees for the Variational Quantum Eigensolver. Version arXiv:2601.16758v2 (CC BY 4.0), distributed by arXiv under the Creative Commons Attribution 4.0 International licence, http://creativecommons.org/licenses/by/4.0/, as recorded in the arXiv licence metadata for this exact version and shown on the abstract page https://arxiv.org/abs/2601.16758v2. Full licence notice: \texttt{licenses/arxiv-2601.16758-CC-BY-4.0.txt}.\par\smallskip

\noindent\textit{Editorial scope.} Counted unit: the article's lemma lemma (source0.tex lines 739--742). The article's complete original proof of it is delivered with it (source0.tex lines 744--778, one \texttt{proof} environment of 35 lines). No mathematics is written, repaired or re-derived: the statement and the proof are the article's own lines, in the article's own notation, and no retained line is substituted or re-flowed.  Retained uncounted as the source-local context the proof needs: lines 686--688; lines 710--712.  Nothing was omitted from any retained span, so the package carries no omission record.  No retained line contains a citation, so the wrapper carries no bibliography and no bibliography entry.  No retained line contains a figure environment, an image-inclusion command or drawing code, so no image is delivered and the package declares no asset dependency.  Editorial formatting only, in addition: an AMS \texttt{amsart} preamble that restates the article's own declarations and macros used by the retained lines, namely the environments \texttt{lemma} on separate counters, together with the standard packages those lines need.  The retained environments print in this document as Lemma 1 in order of appearance, whereas the article prints the same items with the numbering of its own full document.  The complete native source of the article as distributed in the arXiv e-print for this exact version is bundled unchanged as \texttt{sources/193258/source0.tex}.
\begin{align}
    \min_{\theta \in \Theta} \Tr[O \mathcal{E}_\mathbf{\theta}(\rho_0)], \label{perturbedproblem}
\end{align} 

\begin{align}
   \min_{\theta \in \Theta} \Tr[O \mathcal{E}_{L,\theta}(U(\mathbf{\theta}) \rho_0 U(\mathbf{\theta})^\dagger)], \label{perturbedlast}
\end{align}

\begin{lemma}
   Problem \eqref{perturbedproblem} is equivalent to \eqref{perturbedlast} with 
   $p=1-\prod_j (1-p_j)$ and suitably defined operators $\tilde{E}_{j,k}(\theta)$. 
\end{lemma}

\begin{proof}
The proof can be decomposed in two steps.
First, we show that noise acting on any intermediate layer can be seen as a different quantum operation acting on the 
output layer.
Second, we show how to combine nested quantum operations into a single one. 
Combining this two results and using them iteratively for each layer proves the lemma. 

In order to show that mid circuit noises can be represented as a different noise process on the output,
we start defining 
\begin{align}
   \rho_j(\theta)=\mathcal{E}_{j-1,\mathbf{\theta}_{j-1}} \circ \mathcal{E}_{j-2,\mathbf{\theta}_{j-2}} \circ \dots \circ \mathcal{E}_{1,\mathbf{\theta}_1} (\rho_0), 
\end{align}
as the (perturbed) density operator at the $j$-th layer of the circuit. 
Using this, we can define a VQE affected by incoherent noise up to the $j$-th layer as 
\begin{align}
    \min_{\theta \in \Theta} \tilde{\ell} (\theta) = \Tr[O V_{L,j}(\theta) \mathcal{E}_j(\rho_j(\theta))V_{L,j}^\dagger (\theta)],
\end{align}
we can derive the following
\begin{align}
    & \Tr[O V_{L,j}(\theta) \mathcal{E}_j(\rho_j(\theta))V_{L,j}^\dagger (\theta)] \nonumber \\
    & = \Tr[O V_{L,j}(\theta) ((1-p)\rho_j(\theta) + \nonumber p \sum_{k=1}^{m} w_{j,k} E_{j,k} \rho_j(\theta) E^\dagger_{j,k}) V_{L,j}^\dagger (\theta)] \nonumber \\
    & = \Tr[O ((1-p)V_{L,j}(\theta)\rho_j(\theta) V_{L,j}^\dagger (\theta) +  \nonumber \\
    & \quad \quad \quad \quad \quad \quad p \sum_{k=1}^{m} w_{j,k} V_{L,j}(\theta)E_{j,k} \rho_j(\theta) E^\dagger_{j,k}V_{L,j}^\dagger (\theta))]  \nonumber \\
    & = \Tr[O ((1-p)V_{L,j}(\theta)\rho_j(\theta) V_{L,j}^\dagger (\theta) + \nonumber \\
    & \quad \quad \quad \quad \quad \quad p \sum_{k=1}^{m} w_k \tilde{E}_{j,k}(\theta)V_{L,j}(\theta)\rho_j(\theta) V_{L,j}^\dagger (\theta)\tilde{E}^\dagger_{j,k}(\theta))] \nonumber \\
    & = \Tr[O \tilde{\mathcal{E}_\theta}(V_{L,j}(\theta) \rho_j(\theta) V_{L,j}^\dagger (\theta))]    
\end{align}
with $\tilde{E}_{j,k}(\theta) = V_{L,j}(\theta)E_{j,k} V_{L,j}^\dagger (\theta)$. 

For the second step, we use the well-known fact that the composition of quantum operations is itself a quantum operation. 
To derive the error probability $p=1-\prod_j (1-p_j)$ it is sufficient to observe that the probability that the input density operator $\rho$ 
remains unaltered in all the quantum operations is $q=\prod_j (1-p_j)$. It follows immediately that the error probability $p=1-q$.
The Kraus operators can be obtained as all the possible choices of Kraus operators.
By combining the two results and iterating for each layer we get a structure that matches \eqref{perturbedlast}.
\end{proof}

\end{document}
